% =========================================================
% MODELO ABNT - RELATÓRIO TÉCNICO DE SOFTWARE (RTS)
% Instituto Federal do Piauí – Curso Superior de ADS
% Autor - Prof Ricardo Moura Sekeff Budaruiche
% =========================================================

\documentclass[12pt,openright,oneside,a4paper]{article}

% ---------------------------------------------------------
% PACOTES BÁSICOS
\usepackage[T1]{fontenc}
\usepackage[utf8]{inputenc}
\usepackage[brazil]{babel}
\usepackage{times}
\usepackage{setspace}
\usepackage{geometry}
\usepackage{graphicx}
\usepackage{float}
\usepackage{listings}
\usepackage{xcolor}
\usepackage{caption}
\usepackage{booktabs}
\usepackage{hyperref}
\usepackage{amsmath, amssymb}
\usepackage{indentfirst}
\usepackage{lipsum}
\usepackage{enumerate}

\lstset{
    language=Python,            % linguagem padrão (pode trocar)
    basicstyle=\ttfamily\small, % fonte do código
    emph={@app},              % palavras que queremos colorir
    emphstyle=\color{teal}\bfseries,        % estilo delas
    keywordstyle=\color{blue}\bfseries,  % cor das palavras-chave
    backgroundcolor=\color{black!5},
    stringstyle=\color{red},
    commentstyle=\color{gray},
    showstringspaces=false,
    breaklines=true,
    frame=single,               % coloca borda
    numbers=left,               % numeração de linhas
    numberstyle=\tiny\color{gray},
    tabsize=4
}

\renewcommand{\lstlistingname}{Código}
\renewcommand{\lstlistlistingname}{Lista de Códigos}
% ---------------------------------------------------------
% FORMATAÇÃO
% ---------------------------------------------------------
\geometry{a4paper, left=3cm, right=2cm, top=2.5cm, bottom=2.5cm}
\setstretch{1.5}
\captionsetup{labelfont=bf, labelsep=endash}

% ---------------------------------------------------------
% DADOS DO TRABALHO
% ---------------------------------------------------------
\newcommand{\Instituicao}{INSTITUTO FEDERAL DE EDUCAÇÃO, CIÊNCIA E TECNOLOGIA DO PIAUÍ}
\newcommand{\Campus}{Campus Piripiri}
\newcommand{\Curso}{Tecnólogo em Análise e Desenvolvimento de Sistemas}
\newcommand{\Titulo}{MOVY API: API ESCALÁVEL PARA GERENCIAMENTO DE TRANSPORTE UNIVERSITÁRIO INTERMUNICIPAL COM ARQUITETURA MULTITENANCY}
\newcommand{\Autor}{Rauan dos Santos Bandeira}
\newcommand{\Orientador}{Jeferson Soares}
\newcommand{\CoOrientador}{Nome do(a) Coorientador(a) (opcional)}
\newcommand{\Cidade}{Piripiri – PI}
\newcommand{\Ano}{2026}

% ---------------------------------------------------------
\begin{document}

% ================= CAPA =====================
\begin{titlepage}
\begin{center}
\includegraphics[width=10cm]{imagens/Logo-IFPI-Vertical.png} \\[0.3cm]
\textbf{\Instituicao} \\
\textbf{\Campus} \\
\textbf{\Curso} \\[3cm]

\textbf{\Autor} \\[2.5cm]
\textbf{\Titulo} \\[6cm]

\Cidade\\
\Ano
\end{center}
\end{titlepage}

% ================= FOLHA DE ROSTO =====================
\begin{titlepage}
\begin{center}
\textbf{\Autor} \\[3cm]

\textbf{\Titulo} \\[2cm]

\begin{flushright}
Relatório Técnico de Software apresentado ao curso de \Curso{}, como requisito parcial para a conclusão do Trabalho de Conclusão de Curso (TCC). \\
Orientador(a): \Orientador \\
Coorientador(a): \CoOrientador
\end{flushright}

\vfill
\Cidade \\
\Ano
\end{center}
\end{titlepage}

% ================= RESUMO =====================
\newpage
\section*{Resumo}

\textbf{Sumário Executivo}

A Movy API é uma plataforma backend escalável para gerenciamento de transporte universitário intermunicipal, desenvolvida como um Sistema SaaS moderno com suporte a multi-tenancy. O projeto implementa uma arquitetura robusta baseada em Clean Architecture e Domain-Driven Design, utilizando NestJS, TypeScript, Prisma ORM e PostgreSQL. O sistema está 100\% concluído (13 módulos implementados), abrangendo autenticação JWT, gerenciamento de usuários, organizações, roles (RBAC), memberships, motoristas, veículos, viagens (templates e instâncias), automação de agendamento (geração recorrente e auto-cancelamento via \textit{cron}), inscrições de passageiros, planos de assinatura, assinaturas, pagamentos e consumo de plano, com infraestrutura de testes unitários composta por 458 testes em 57 suites.

\vspace{0.5cm}

\noindent\textbf{Destaques Técnicos}:
\begin{itemize}
\item Multi-tenancy estruturado com isolamento lógico de dados por organização
\item Autenticação JWT com Bcrypt (10 salt rounds) e refresh token
\item RBAC completo com guards (RolesGuard, TenantFilterGuard, DevGuard)
\item Value Objects com validações encapsuladas (CNPJ, Email, Slug, CNH, Placa, Money)
\item Soft-delete em todas as entidades críticas
\item Docker Compose ready com auto-seed de dados e migrations automáticas
\item Swagger integrado com documentação automática de endpoints
\item Migrações versionadas (28 migrações executadas)
\item Rate limiting global (60 req/min via \texttt{@nestjs/throttler})
\item Proteção contra IDOR (OWASP A01) em todos os recursos por ID
\item Transações atômicas via UnitOfWork + DbContext (AsyncLocalStorage)
\item Controle de limites de plano via \texttt{PlanLimitService} (veículos, motoristas, viagens/mês)
\item Expiração lazy de assinaturas (sem cron job)
\item Testes unitários: 458 testes, 57 suites, padrão AAA com injeção manual
\item CI/CD configurado via GitHub Actions
\end{itemize}

\vspace{0.5cm}

\noindent\textbf{Descrição Técnica Detalhada}\vspace{0.3cm}

A Movy API é uma solução backend escalável, desenvolvida como um Sistema de Software como Serviço (SaaS), destinada ao gerenciamento de transporte universitário intermunicipal. O sistema implementa uma arquitetura de multi-tenancy que permite que múltiplas organizações de transporte utilizem simultaneamente a plataforma, com isolamento lógico de dados. Utilizando princípios de Domain-Driven Design (DDD) e Clean Architecture, a API foi desenvolvida em NestJS com TypeScript, Prisma ORM e PostgreSQL, garantindo escalabilidade, segurança e manutenibilidade. O projeto implementa autenticação JWT com controle de acesso baseado em roles (RBAC), permitindo diferentes níveis de permissão para administradores e motoristas. Foram implementados 13 módulos funcionais — usuários, autenticação, organizações, roles, memberships, motoristas, veículos, viagens (templates e instâncias), inscrições (bookings), planos, assinaturas e pagamentos — totalizando 100\% do escopo de módulos previsto. A arquitetura proposta segue padrões estabelecidos de engenharia de software, facilitando integrações com aplicações front-end e sistemas externos.

\textbf{Palavras-chave:} Transporte universitário, SaaS, API REST, Multi-tenancy, NestJS, DDD, Clean Architecture, Autenticação JWT.

\newpage
\section*{Abstract}
The Movy API is a scalable backend solution developed as a Software as a Service (SaaS) system for managing intermunicipal university transportation. The system implements a multi-tenancy architecture that allows multiple transportation organizations to simultaneously use the platform with logical data isolation. Using Domain-Driven Design (DDD) and Clean Architecture principles, the API was developed in NestJS with TypeScript, Prisma ORM, and PostgreSQL, ensuring scalability, security, and maintainability. The project implements JWT authentication with role-based access control (RBAC), allowing different permission levels for administrators and drivers. Thirteen functional modules have been fully implemented — users, authentication, organizations, roles, memberships, drivers, vehicles, trips (templates and instances), passenger bookings, plans, subscriptions, and payments — totaling 100\% of the planned module scope. The proposed architecture follows established software engineering patterns, facilitating integration with front-end applications and external systems.

\textbf{Keywords\:} University transportation, SaaS, REST API, Multi-tenancy, NestJS, DDD, Clean Architecture, JWT Authentication.

% ================= SUMÁRIO =====================
\newpage
\tableofcontents
% \lstlistoflistings % Para exibir a listá de listagens de códigos
\newpage

% ================= DESENVOLVIMENTO =====================
\section{Introdução}
\label{sec:introducao}
O ensino superior presencial no Brasil apresenta uma expressiva dimensão e mantém milhões de estudantes matriculados ao longo dos últimos anos. Dados do Censo da Educação Superior, divulgados pelo Instituto Nacional de Estudos e Pesquisas Educacionais Anísio Teixeira (INEP), indicam que, em 2010, o país registrava 5.476.813 matrículas presenciais, número que se manteve elevado, alcançando 5.037.875 matrículas em 2024, distribuídas em 2.444 instituições de ensino superior presenciais \cite{inepCensoSuperior}. Esses dados evidenciam que uma parcela significativa dos universitários depende da frequência física às instituições para a realização de suas atividades acadêmicas.

Essa realidade é reforçada por informações do Censo Demográfico de 2022, realizado pelo Instituto Brasileiro de Geografia e Estatística (IBGE), que apontam que 27,3\% dos estudantes de graduação cursam o ensino superior em município diferente daquele em que residem \cite{ibgeDeslocamentoEstudo}. Esse cenário evidencia que o deslocamento intermunicipal faz parte da rotina de uma parcela significativa dos universitários, especialmente em regiões onde a oferta de cursos presenciais está concentrada em cidades-polo.

Apesar da importância do transporte universitário intermunicipal para a permanência desses estudantes no ensino superior, a organização desse serviço ainda ocorre, em muitas localidades, de forma manual e pouco estruturada. O agendamento das viagens e o controle de presença são frequentemente realizados por meio de listas informais em aplicativos de mensagens, como grupos de WhatsApp, nos quais os alunos registram seus nomes para confirmar a participação. Esse modelo de gestão está sujeito a falhas, como duplicidade de registros, esquecimentos, dificuldades de comunicação e ausência de controle financeiro adequado.

O cenário foi observado diretamente na prática, durante a utilização diária do serviço em uma rota específica.  Constatou-se que a inexistência de um sistema informatizado compromete a visualização clara dos passageiros confirmados, dificulta o gerenciamento dos pagamentos e aumenta o risco de decisões equivocadas quanto à viabilidade das viagens, uma vez que, em muitos casos, há um número mínimo de passageiros necessário para que o deslocamento ocorra. Dessa forma, problemas como atrasos, prejuízos financeiros e desorganização tornam-se recorrentes.

Diante desse contexto, evidencia-se a necessidade de uma solução tecnológica capaz de automatizar o processo de agendamento de viagens, o registro de passageiros e o controle financeiro do transporte universitário intermunicipal, proporcionando maior organização, confiabilidade e eficiência tanto para os estudantes quanto para os responsáveis pelo serviço.


\subsection{Objetivos}

\subsubsection{Objetivo Geral}
Desenvolver uma API (Interface de Programação de Aplicativos) escalável para o gerenciamento do transporte universitário intermunicipal, com o objetivo de automatizar o agendamento de viagens, o registro de passageiros e o controle financeiro. A solução visa proporcionar maior organização, transparência e eficiência para diversos clientes (estudantes, responsáveis pelo transporte e possíveis aplicações front-end).

\subsubsection{Objetivos Específicos}
\begin{itemize}
\item Analisar o funcionamento atual do transporte universitário intermunicipal, identificando problemas e limitações dos processos manuais adotados;
\item Levantar e documentar os requisitos funcionais e não funcionais do sistema, com base nas necessidades dos estudantes e dos responsáveis pelo transporte;
\item Projetar a arquitetura da API, definindo padrões de comunicação, organização das camadas e estrutura do banco de dados;
\item Modelar os principais artefatos do sistema, incluindo casos de uso e modelo entidade-relacionamento;
\item Implementar a API REST utilizando tecnologias adequadas ao desenvolvimento de sistemas back-end;
\item Realizar testes funcionais e unitários para validar o correto funcionamento das principais funcionalidades implementadas;
\item Avaliar os resultados obtidos com a solução desenvolvida, considerando aspectos de organização, confiabilidade e facilidade de uso.
\end{itemize}


% ------------------------------
\section{Trabalhos Relacionados}

Diversos estudos têm abordado o uso de sistemas de informação no gerenciamento de transporte estudantil e universitário, com o objetivo de otimizar processos, reduzir falhas operacionais e melhorar a organização das viagens.

O trabalho PickBus, proposto por Pretto et al. \cite{pickbus2021}, apresenta uma solução multiplataforma (web e mobile) voltada ao gerenciamento de transporte fretado de passageiros, com foco em eventos e deslocamentos coletivos. O sistema permite a criação e o gerenciamento de viagens, controle de passageiros, organização de pagamentos e centralização da comunicação entre organizadores e usuários. O estudo fundamenta-se em conceitos de mobilidade urbana e empreendedorismo, demonstrando a viabilidade da solução por meio de pesquisa de mercado e testes com usuários reais. Os resultados indicam boa aceitação do sistema e evidenciam a carência de ferramentas digitais específicas para o gerenciamento de transporte fretado.


O trabalho intitulado "Sistema para Gerenciamento de Transportes", desenvolvido por Anderson et al. \cite{anderson2020}, apresenta uma solução voltada à modernização do processo de solicitação e gerenciamento de transportes no contexto de uma instituição acadêmica. A proposta surge da necessidade de substituir métodos manuais baseados em troca de e-mails, os quais são suscetíveis a falhas de comunicação e perda de informações. O sistema permite o cadastro e a gestão de veículos, motoristas e setores, bem como a solicitação de viagens, verificação de disponibilidade, fluxo de aprovação hierárquico e alocação de recursos. Além disso, a solução contempla funcionalidades de monitoramento por meio de painéis e geração de relatórios, contribuindo para maior transparência e apoio à tomada de decisão. O estudo evidencia que a informatização do processo de gestão de transportes institucionais promove maior organização, eficiência operacional e confiabilidade das informações.


Outro trabalho relevante é o de Reis et al. \cite{reis2021transporte}, que apresenta a modelagem de um software voltado ao gerenciamento do transporte escolar público municipal. A proposta tem como foco o levantamento de requisitos e a utilização de técnicas de Engenharia de Software e UML para organizar informações relacionadas a alunos, rotas, veículos, motoristas e monitores. O estudo evidencia que a ausência de sistemas informatizados compromete a eficiência operacional e a tomada de decisão, especialmente em contextos com grande volume de usuários e necessidade de relatórios precisos. Diferentemente do presente projeto, que se concentra no transporte universitário intermunicipal e na oferta de uma API escalável para integração com diferentes aplicações cliente, o trabalho de Reis et al. (2021) enfatiza a modelagem conceitual e o contexto do ensino básico municipal, reforçando, contudo, a importância da informatização como fator essencial para a organização e confiabilidade dos serviços de transporte educacional.

O sistema proposto neste trabalho diferencia-se das soluções analisadas por abordar de forma específica o contexto do transporte universitário intermunicipal, um cenário ainda pouco explorado na literatura, sobretudo no que se refere à integração entre organização operacional e controle financeiro. Enquanto os trabalhos relacionados concentram-se no transporte escolar, no fretamento eventual para eventos ou na gestão institucional de veículos, a presente proposta foca nas particularidades do deslocamento diário e recorrente de estudantes universitários entre municípios, considerando suas demandas operacionais e administrativas.

Além da automatização do agendamento de viagens, do gerenciamento da lista de passageiros e da centralização da comunicação entre usuários e responsáveis pelo transporte, o principal diferencial do sistema está na inclusão de um módulo de controle financeiro integrado à gestão das viagens. O sistema permite a estimativa automática da arrecadação por viagem com base no número de passageiros confirmados, possibilitando ao responsável avaliar rapidamente a viabilidade financeira de cada deslocamento.

Adicionalmente, ao ser concebido como uma API escalável com modelo SaaS multi-tenant, o projeto amplia sua aplicabilidade, permitindo a integração com diferentes sistemas clientes, como aplicações web ou móveis, o que não é observado nos trabalhos analisados.


% ------------------------------
\section{Tecnologias Implementadas}

Os requisitos de escalabilidade, segurança e manutenibilidade do projeto conduziram à seleção de uma stack tecnológica robusta e contemporânea. A adoção de NestJS foi fundamentada na necessidade de uma arquitetura estruturada, com suporte nativo a padrões como injeção de dependência, módulos, guards e interceptadores globais, essenciais para implementar Clean Architecture e RBAC de forma consistente.

\subsection{Stack Tecnológico Adotado}

\begin{table}[H]
\centering
\caption{Tecnologias implementadas no desenvolvimento da Movy API}
\label{tab:tecnologias_implementadas}
\begin{tabular}{lll}
\toprule
\textbf{Tecnologia} & \textbf{Versão} & \textbf{Descrição} \\
\midrule
Node.js & v18+ & Ambiente de execução JavaScript/TypeScript \\
TypeScript & 5.3+ & Linguagem tipada para desenvolvimento seguro \\
NestJS & v11.0+ & Framework backend modular com DI nativo\\
Prisma ORM & v7.5+ & ORM com tipagem forte e migrations automáticas \\
PostgreSQL & 15+ & SGBD relacional para persistência de dados \\
JWT (NestJS/Passport) & 11.0+ & Implementação de autenticação com tokens \\
Bcrypt & 6.0+ & Hash seguro de senhas com salt rounds \\
Docker & Latest & Conteinerização e orquestração de ambientes \\
Jest & 30+ & Framework de testes unitários e de integração \\
ESLint & 9+ & Análise estática de código TypeScript \\
Swagger & 11.2+ & Documentação interativa de endpoints REST \\
@nestjs/schedule & 6.1+ & Agendamento de \textit{jobs} (cron) de viagens \\
@nestjs/throttler & 6.5+ & \textit{Rate limiting} global (60 req/min) \\
\bottomrule
\end{tabular}
\end{table}

\subsection{Justificativa das Escolhas Tecnológicas}

\subsubsection{NestJS}
NestJS fornece:
\begin{itemize}
\item \textbf{Injeção de Dependência Nativa}: Essencial para Clean Architecture e padrão Repositório.
\item \textbf{Módulos Estruturados}: Permite organização clara de contextos delimitados (bounded contexts) do DDD.
\item \textbf{Guards e Interceptadores Globais}: Implementação seamless de RBAC e middleware de logging.
\item \textbf{Ecosystem Maduro}: Suporte a Swagger, JWT, Passport e outras bibliotecas essenciais.
\item \textbf{Escalabilidade}: Pronto para microserviços e processamento em background com agenda de jobs.
\end{itemize}

\subsubsection{PostgreSQL e Prisma ORM}
\begin{itemize}
\item \textbf{PostgreSQL}: Banco de dados relacional confiável com suporte a ACID, transações complexas e tipos de dados avançados.
\item \textbf{Prisma}: Oferece:
\begin{enumerate}
\item Tipagem forte: Geração automática de tipos TypeScript a partir do schema.
\item Migrations versionadas: Histórico completo de mudanças na estrutura do banco.
\item Query builder seguro contra SQL injection.
\item Suporte nativo a soft-delete e operações transacionais.
\item Seed scripts para inicialização de dados.
\end{enumerate}
\end{itemize}

\subsubsection{JWT com Passport.js}
Implementação de autenticação stateless, adequada para APIs REST e arquiteturas de microserviços. O refresh token permite renovação de credenciais sem re-autenticação. O payload do JWT é enriquecido no momento do login com todos os dados necessários (\texttt{userId}, \texttt{organizationId}, \texttt{role}, \texttt{isDev}), eliminando consultas ao banco a cada requisição autenticada.

\subsubsection{Docker para Containerização}
\begin{itemize}
\item Padronização de ambientes (desenvolvimento, teste, produção).
\item Facilita deploy em plataformas cloud (AWS, GCP, Azure, Render, etc.).
\item Docker Compose para orquestração local com PostgreSQL + API sincronizados.
\end{itemize}

\subsubsection{Jest para Testes}
Framework padrão da comunidade Node/NestJS, com suporte a testes unitários, de integração e cobertura de código. Integração simplificada com TypeScript via ts-jest.


% ------------------------------
\section{Modelagem do Projeto}

\subsection{Levantamento de Requisitos}
Os requisitos do sistema foram definidos a partir da observação do funcionamento atual do transporte universitário intermunicipal, bem como das necessidades relatadas por estudantes, motoristas e responsáveis pela gestão das viagens. Para melhor organização, os requisitos foram classificados em funcionais, não funcionais e de domínio. \\
\textbf{Requisitos Funcionais}:
\begin{itemize}
\item RF01 – Permitir o cadastro de usuários, contendo informações pessoais e dados de identificação e possível vinculação à organização responsável.
\item RF02 – Permitir o cadastro de organizações e administradores responsáveis pela gestão das viagens.
\item RF03 – Permitir o cadastro de motoristas vinculados às viagens.
\item RF04 – Permitir o agendamento de viagens com data, horário específico, turno e destino.
\item RF05 – Permitir o cancelamento de viagens pela organização ou administrador responsável.
\item RF06 – Permitir a definição de turnos das viagens (manhã, tarde, noite ou outros).
\item RF07 – Permitir que usuários se inscrevam em viagens disponíveis.
\item RF08 – Permitir o cancelamento da inscrição do usuário dentro de um prazo definido.
\item RF09 – Permitir a escolha do tipo de trajeto (ida, volta ou ida e volta).
\item RF10 – Permitir a indicação da forma de pagamento associada à viagem (por exemplo, PIX ou dinheiro).
\item RF11 – Permitir a visualização da lista de passageiros inscritos em cada viagem.
\item RF12 – Permitir a confirmação de presença do passageiro no momento do embarque.
\item RF13 – Permitir que o motorista visualize a lista de passageiros da viagem.
\item RF14 – Permitir que o motorista marque a viagem como iniciada ou concluída.
\item RF15 – Permitir a geração de relatórios das viagens realizadas (possível redução de escopo atual).
\item RF16 – Permitir que a organização visualize o histórico de viagens (possível redução de escopo atual).
\item RF17 – Permitir o envio de notificações para lembrar os usuários sobre inscrições em viagens (possível redução de escopo atual).
\item RF18 – Permitir o envio de notificações de confirmação ou cancelamento (possível redução de escopo atual).
\item RF19 – Calcular a estimativa financeira de arrecadação por viagem.
\item RF20 – Permitir o cadastro de meios de transporte (ônibus, vans, entre outros).
\item RF21 – Permitir a vinculação de meios de transporte e motoristas às viagens.
\item RF22 – Impedir inscrições acima da capacidade máxima definida para o transporte.
\item RF23 – Permitir o cadastro de pontos de embarque pela organização (possível redução de escopo atual).
\item RF24 – Permitir que o usuário selecione um ponto de embarque no momento da inscrição (possível simplificação no escopo atual).
\item RF25 – Permitir que o motorista visualize os pontos de embarque associados às inscrições da viagem (possível redução de escopo atual)
\item RF26 - Registro de Logs de auditoria.

\end{itemize}

\textbf{Requisitos Não Funcionais}:
\begin{itemize}
\item RNF01 (Usabilidade da API) – A API deve seguir o padrão RESTful, utilizando métodos HTTP adequados e códigos de status padronizados para garantir uma integração intuitiva por desenvolvedores e futuros front-ends.
\item RNF02 (Performance) – O sistema deve apresentar tempo de resposta (latência) inferior a 2 segundos para requisições de consulta, garantindo fluidez no processamento de dados, conforme recomendado por Pressman \cite{pressman2016}.

\item RNF03 (Disponibilidade) – O sistema deve possuir disponibilidade mínima de 95\%, garantindo que os serviços de transporte não sejam interrompidos em horários críticos de operação.

\item RNF04 (Segurança e Privacidade) – O sistema deve implementar autenticação via JWT (JSON Web Token) e controle de acesso baseado em funções (RBAC), garantindo a proteção de dados sensíveis em conformidade com os princípios da LGPD.

\item RNF05 (Escalabilidade e Multi-tenancy) – A arquitetura deve permitir escalabilidade horizontal e isolamento lógico de dados por organização, permitindo que o sistema funcione como um modelo SaaS (Software as a Service).

\item RNF06 (Manutenibilidade) – O código deve ser desenvolvido utilizando o framework NestJS, seguindo princípios de Clean Architecture e SOLID, para facilitar manutenções e evoluções futuras.

\item RNF07 (Interoperabilidade) – O sistema deve ser independente de plataforma, comunicando-se exclusivamente via JSON, permitindo o consumo por aplicativos móveis ou sistemas web.

\item RNF08 (Integridade de Dados) – O sistema deve garantir a consistência das informações através de transações ACID no banco de dados, evitando duplicidade em inscrições e garantindo a precisão do controle de vagas.

\item RNF09 (Persistência e Confiabilidade) – Os dados devem ser armazenados em um sistema de gerenciamento de banco de dados relacional (PostgreSQL), garantindo a persistência a longo prazo e a segurança contra perda de informações.
\end{itemize}

\subsection{Casos de Uso}
\begin{figure}[H]
\centering
\includegraphics[width=0.6\textwidth]{imagens/usecases_diagram.png}
\caption{Diagrama de caso de uso.}
\label{fig:sapo}
\end{figure}

\subsection{Diagrama de Classes}

O modelo de domínio foi modelado segundo os princípios do \textit{Domain-Driven Design} (DDD): cada entidade é um agregado com métodos de fábrica estáticos (\texttt{create()}, que valida invariantes, e \texttt{restore()}, que reidrata a partir da persistência sem revalidar), e os atributos de regra de negócio são encapsulados em \textit{Value Objects} imutáveis (por exemplo, \texttt{Email}, \texttt{Cnpj}, \texttt{Cnh}, \texttt{Plate}, \texttt{Money}). Para preservar a legibilidade em formato impresso, o diagrama é apresentado em um \textbf{mapa panorâmico} (Figura~\ref{fig:classes-geral}), que exibe apenas as entidades e suas associações, seguido de figuras de detalhe organizadas por \textit{bounded context} (identidade e acesso, frota, viagens, reservas e pagamentos, e \textit{billing}), nas quais constam atributos, métodos e \textit{Value Objects}. O detalhamento completo encontra-se no arquivo-fonte \texttt{docs/DIAGRAMA\_CLASSES.md}, do qual as figuras são exportadas.

\begin{figure}[H]
\centering
% TODO: exportar de docs/DIAGRAMA_CLASSES.md (Fig. 1) e salvar em imagens/
% \includegraphics[width=\textwidth]{imagens/diagrama_classes_geral.png}
\fbox{\parbox[c][4cm][c]{0.9\textwidth}{\centering\textit{[Mapa geral de entidades --- exportar da Fig.~1 de \texttt{docs/DIAGRAMA\_CLASSES.md}]}}}
\caption{Mapa geral de entidades do modelo de domínio (Movy API).}
\label{fig:classes-geral}
\end{figure}

As entidades agrupam-se em cinco contextos: \textbf{(i) Identidade e Acesso} --- \texttt{User}, \texttt{Organization}, \texttt{Membership} (tabela-pivô que materializa o N:N entre usuário e organização, carregando o \texttt{Role} e servindo de base ao RBAC) e os \textit{tokens} de redefinição de senha e verificação de e-mail; \textbf{(ii) Frota} --- \texttt{Driver} (entidade global, vinculada a organizações via \texttt{Membership}) e \texttt{Vehicle}; \textbf{(iii) Viagens} --- \texttt{TripTemplate} (modelo recorrente) que gera \texttt{TripInstance} (execução datada, com máquina de estados) e \texttt{TripSchedulingConfig}; \textbf{(iv) Reservas e Pagamentos} --- \texttt{Booking} (inscrição) e \texttt{Payment} (1:1 com a inscrição); e \textbf{(v) Billing} --- \texttt{Plan} e \texttt{Subscription}, com cota de uso aplicada pelo \texttt{PlanLimitService}.

\subsection{Arquitetura do Sistema}

A arquitetura foi estruturada seguindo os princípios da Clean Architecture e Domain-Driven Design (DDD), organizando o sistema em camadas bem definidas que isolam responsabilidades e permitem independência de frameworks externos.

\subsubsection{Camadas Arquiteturais}

\begin{enumerate}
\item \textbf{Camada de Apresentação (Presentation)}: Controladores REST que orquestram requisições HTTP, validação de DTOs, autenticação JWT e retorno de respostas formatadas.
\begin{itemize}
    \item Controllers: Mapeiam rotas HTTP para casos de uso.
    \item DTOs: Data Transfer Objects com validação via class-validator.
    \item Presenters: Formatam respostas de domínio para o cliente.
    \item Mappers: Convertem entidades de domínio em respostas HTTP.
\end{itemize}

\item \textbf{Camada de Aplicação (Application)}: Orquestra regras de negócio através de Use Cases.
\begin{itemize}
    \item Use Cases: Implementam a lógica de negócio de forma testável e independente.
    \item Aplicação: Coordena fluxo entre domínio e infraestrutura.
    \item Exceptions: Erros de aplicação específicos (ex: \textit{MembershipAlreadyExistsError}).
\end{itemize}

\item \textbf{Camada de Domínio (Domain)}: Contém regras de negócio centrais, completamente independentes de frameworks.
\begin{itemize}
    \item Entities: Objetos com identidade única (ex: User, Organization).
    \item Value Objects: Encapsulam validações de dados (ex: Email, CNPJ, Slug, Money).
    \item Interfaces: Abstraem contratos de repositórios e serviços.
    \item Erros de Domínio: Exceções específicas do negócio.
\end{itemize}

\item \textbf{Camada de Infraestrutura (Infrastructure)}: Implementação de detalhes técnicos.
\begin{itemize}
    \item Repositories: Implementação do padrão Repositório com Prisma.
    \item Database: Configuração do Prisma e conexão com PostgreSQL.
    \item Mappers: Conversão entre entidades de domínio e modelos de persistência.
    \item Providers: Serviços específicos (ex: Bcrypt para hash de senhas).
\end{itemize}

\item \textbf{Camada Compartilhada (Shared)}: Componentes reutilizáveis.
\begin{itemize}
    \item Domain Errors e Value Objects compartilhados.
    \item Guards (JWT authentication, RolesGuard, TenantFilterGuard, DevGuard).
    \item Interceptadores de logging.
    \item Filters de exceção global.
    \item UnitOfWork e DbContext para transações atômicas.
    \item PrismaModule para injeção da conexão em toda a aplicação.
\end{itemize}
\end{enumerate}

\subsubsection{Estrutura de Módulos}

Cada módulo de negócio segue a mesma estrutura hierárquica:

\begin{lstlisting}[caption={Estrutura de um módulo (exemplo: usuários)}, label={cod:modulo_estrutura}]
src/modules/user/
  application/
    dto/
      - create-user.dto.ts
      - update-user.dto.ts
      - user-response.dto.ts
    use-cases/
      - create-user.use-case.ts
      - find-all-users.use-case.ts
      - find-user-by-id.use-case.ts
  domain/
    entities/
      - user.entity.ts
    errors/
      - user.errors.ts
    interfaces/
      - user.repository.ts
  infrastructure/
    db/
      mappers/
        - user.mapper.ts
      repositories/
        - prisma-user.repository.ts
  presentation/
    controllers/
      - user.controller.ts
    mappers/
      - user.mapper.ts
  user.module.ts
\end{lstlisting}

\subsubsection{Padrões de Design Implementados}

\begin{enumerate}
\item \textbf{Padrão Repositório}: Abstração da persistência de dados.
\item \textbf{Injeção de Dependência}: Gerenciada nativamente pelo NestJS.
\item \textbf{Value Objects}: Encapsulam validações de valores complexos.
\item \textbf{Use Cases}: Cada operação de negócio é uma classe testável independente.
\item \textbf{Guards}: Proteção de rotas com autenticação JWT.
\item \textbf{Interceptadores}: Logging centrado e manipulação de respostas.
\item \textbf{Filters de Exceção}: Tratamento global de erros com formatação padronizada.
\item \textbf{Soft Delete}: Marcação de exclusão via status/timestamp.
\item \textbf{UnitOfWork}: Transações atômicas via AsyncLocalStorage sem expor Prisma à camada de aplicação.
\end{enumerate}

\subsubsection{Pipeline de Autenticação e Autorização}

O pipeline de segurança é composto por quatro guards executados em sequência:

\begin{lstlisting}[caption={Pipeline de autorização}, label={cod:auth_pipeline}]
Request
  -> JwtAuthGuard      # valida JWT; constroi TenantContext -> req.context
  -> RolesGuard        # verifica @Roles(); isDev ignora
  -> TenantFilterGuard # :organizationId param == req.context.organizationId
  -> DevGuard          # apenas isDev=true passa em rotas @Dev()
  -> Controller
\end{lstlisting}

O \texttt{JwtAuthGuard} enriquece o payload JWT no momento do login com todos os dados necessários (\texttt{userId}, \texttt{organizationId}, \texttt{role}, \texttt{isDev}), eliminando consultas ao banco a cada requisição. Usuários listados em \texttt{DEV\_EMAILS} recebem \texttt{isDev=true} no token e ignoram validações de tenant e role.

\subsection{Diagrama Entidade-Relacionamento}

O Diagrama Entidade-Relacionamento (DER) foi elaborado com o objetivo de representar a estrutura lógica do banco de dados do sistema, evidenciando as principais entidades, seus atributos e os relacionamentos existentes entre elas. O modelo foi construído considerando as regras de negócio do transporte universitário intermunicipal, garantindo coerência, integridade dos dados e suporte às funcionalidades propostas pela API.

A modelagem reflete um design de multi-tenancy, onde cada organização possui isolamento lógico de dados através do campo \textit{organizationId}. As principais entidades implementadas incluem: Usuário, Organização, Roles, Memberships, Motoristas, Veículos, Templates de Viagem, Instâncias de Viagem, Inscrições, Pagamentos, Planos e Assinaturas, além de registros de Auditoria.

\begin{figure}[H]
\centering
\includegraphics[width=\textwidth]{imagens/Subscription Payment-2026-04-06-214838.png}
\caption{Diagrama Entidade-Relacionamento do sistema (Movy API).}
\label{fig:der}
\end{figure}

\subsubsection{Entidades Implementadas}

\paragraph{\textbf{Role (Papéis de Acesso)}}\mbox{}\\
Armazena os diferentes papéis (roles) disponíveis no sistema para controle de acesso baseado em funções (RBAC).
\begin{itemize}
    \item \textbf{id} -- Identificador único (autoincrement)
    \item \textbf{name} -- Nome do papel (ADMIN, DRIVER)
    \item \textbf{createdAt} -- Timestamp de criação
    \item \textbf{updatedAt} -- Timestamp de atualização
\end{itemize}

\paragraph{\textbf{User (Usuários)}}\mbox{}\\
Armazena informações de todos os usuários do sistema.
\begin{itemize}
    \item \textbf{id} -- Identificador único (UUID)
    \item \textbf{name} -- Nome completo do usuário
    \item \textbf{email} -- Endereço de e-mail único
    \item \textbf{passwordHash} -- Senha criptografada com Bcrypt
    \item \textbf{telephone} -- Telefone para contato
    \item \textbf{status} -- Status do usuário (ACTIVE, INACTIVE) - soft delete
    \item \textbf{createdAt} -- Timestamp de criação
    \item \textbf{updatedAt} -- Timestamp de atualização
\end{itemize}

\paragraph{\textbf{Organization (Organizações)}}\mbox{}\\
Armazena dados das organizações de transporte cadastradas no sistema.
\begin{itemize}
    \item \textbf{id} -- Identificador único (UUID)
    \item \textbf{name} -- Nome da organização
    \item \textbf{cnpj} -- CNPJ da organização (único, com validação de dígitos verificadores)
    \item \textbf{email} -- E-mail institucional único
    \item \textbf{telephone} -- Telefone para contato
    \item \textbf{address} -- Endereço completo
    \item \textbf{slug} -- Identificador URL-friendly único
    \item \textbf{status} -- Status (ACTIVE, INACTIVE) - soft delete
    \item \textbf{createdAt} -- Timestamp de criação
    \item \textbf{updatedAt} -- Timestamp de atualização
\end{itemize}

\paragraph{\textbf{Plan (Planos de Assinatura)}}\mbox{}\\
Define os planos (tiers) de serviço disponíveis (FREE, BASIC, PRO, PREMIUM).
\begin{itemize}
    \item \textbf{id} -- Identificador único (autoincrement)
    \item \textbf{name} -- Nome do plano (FREE, BASIC, PRO, PREMIUM)
    \item \textbf{price} -- Preço (Decimal 10.2)
    \item \textbf{description} -- Descrição do plano
    \item \textbf{maxVehicles} -- Quantidade máxima de veículos
    \item \textbf{maxDrivers} -- Quantidade máxima de motoristas
    \item \textbf{maxMonthlyTrips} -- Quantidade máxima de viagens por mês
    \item \textbf{durationDays} -- Duração da assinatura em dias (padrão: 30)
    \item \textbf{isActive} -- Se o plano está disponível
    \item \textbf{createdAt, updatedAt} -- Timestamps
\end{itemize}

\paragraph{\textbf{Subscription (Assinaturas)}}\mbox{}\\
Registra as assinaturas das organizações nos planos.
\begin{itemize}
    \item \textbf{id} -- Identificador único (UUID)
    \item \textbf{organizationId} -- FK: Organização
    \item \textbf{planId} -- FK: Plano
    \item \textbf{status} -- Status da assinatura (ACTIVE, PAST\_DUE, CANCELED)
    \item \textbf{startDate} -- Data de início
    \item \textbf{expiresAt} -- Data de expiração (startDate + plan.durationDays)
    \item \textbf{createdAt, updatedAt} -- Timestamps
\end{itemize}

\paragraph{\textbf{OrganizationMembership (Associações Usuário-Role-Organização)}}\mbox{}\\
Tabela de junção que associa usuários com roles em organizações específicas, formando a base do RBAC.
\begin{itemize}
    \item \textbf{userId, roleId, organizationId} -- Chave composta (PK)
    \item \textbf{assignedAt} -- Data de atribuição do papel
    \item \textbf{removedAt} -- Data de remoção (NULL para memberships ativos) - soft delete
    \item Relacionamentos: User (cascade), Role (cascade), Organization (cascade)
\end{itemize}

\paragraph{\textbf{Driver (Motoristas)}}\mbox{}\\
Extensão da entidade User para motoristas. Motoristas são entidades globais, vinculados a organizações exclusivamente via \texttt{OrganizationMembership}, permitindo que um usuário seja motorista em múltiplas organizações.
\begin{itemize}
    \item \textbf{userId} -- PK e FK: Usuário
    \item \textbf{cnh} -- Número da CNH (único)
    \item \textbf{cnhCategories} -- Categorias da CNH (coleção; um motorista pode ter A+B etc.)
    \item \textbf{cnhExpiresAt} -- Data de validade da CNH
    \item \textbf{status} -- Status (ACTIVE, INACTIVE, SUSPENDED)
    \item \textbf{createdAt, updatedAt} -- Timestamps
\end{itemize}

\paragraph{\textbf{Vehicle (Veículos)}}\mbox{}\\
Armazena informações dos veículos da frota.
\begin{itemize}
    \item \textbf{id} -- Identificador único (UUID)
    \item \textbf{plate} -- Placa do veículo (único, formatos ABC1234 e Mercosul ABC1D23)
    \item \textbf{model} -- Modelo do veículo
    \item \textbf{type} -- Tipo (VAN, BUS, MINIBUS, CAR)
    \item \textbf{maxCapacity} -- Capacidade máxima de passageiros
    \item \textbf{status} -- Status (ACTIVE, INACTIVE)
    \item \textbf{organizationId} -- FK: Organização (multi-tenancy)
    \item \textbf{createdAt, updatedAt} -- Timestamps
\end{itemize}

\paragraph{\textbf{TripTemplate (Template de Viagens Recorrentes)}}\mbox{}\\
Define template para viagens recorrentes (rotas com horários e frequência).
\begin{itemize}
    \item \textbf{id} -- Identificador único (UUID)
    \item \textbf{departurePoint} -- Local de partida
    \item \textbf{destination} -- Destino
    \item \textbf{departureTimeOfDay, arrivalTimeOfDay} -- Hora do dia (HH:mm, UTC) de partida/chegada
    \item \textbf{priceOneWay, priceReturn, priceRoundTrip} -- Preços por tipo de trajeto
    \item \textbf{isRecurring} -- Se é recorrente
    \item \textbf{frequency} -- Array de dias da semana (DayOfWeek enum)
    \item \textbf{shift} -- Turno (MORNING, AFTERNOON, EVENING)
    \item \textbf{stops} -- Array de paradas
    \item \textbf{defaultCapacity} -- Capacidade padrão propagada às instâncias geradas
    \item \textbf{autoCancelEnabled, minRevenue, autoCancelOffset} -- Configuração de auto-cancelamento
    \item \textbf{status} -- Status (ACTIVE, INACTIVE)
    \item \textbf{organizationId} -- FK: Organização
    \item \textbf{isPublic} -- Se está disponível para inscrição
    \item \textbf{createdAt, updatedAt} -- Timestamps
\end{itemize}

\paragraph{\textbf{TripInstance (Instâncias de Viagem)}}\mbox{}\\
Representa uma viagem específica (execução de um template).
\begin{itemize}
    \item \textbf{id} -- Identificador único (UUID)
    \item \textbf{driverId} -- FK: Motorista (nullable)
    \item \textbf{vehicleId} -- FK: Veículo (nullable)
    \item \textbf{totalCapacity} -- Capacidade total da viagem
    \item \textbf{tripStatus} -- Status (DRAFT, SCHEDULED, CONFIRMED, IN\_PROGRESS, FINISHED, CANCELED)
    \item \textbf{tripTemplateId} -- FK: Template
    \item \textbf{departureTime} -- Horário de partida
    \item \textbf{arrivalEstimate} -- Previsão de chegada
    \item \textbf{organizationId} -- FK: Organização (multi-tenancy)
    \item \textbf{createdAt, updatedAt} -- Timestamps
\end{itemize}

\paragraph{\textbf{Booking (Inscrições)}}\mbox{}\\
Registra a inscrição de um usuário em uma viagem específica.
\begin{itemize}
    \item \textbf{id} -- Identificador único (UUID)
    \item \textbf{enrollmentDate} -- Data da inscrição
    \item \textbf{status} -- Status (ACTIVE, INACTIVE) - soft delete
    \item \textbf{presenceConfirmed} -- Se a presença foi confirmada
    \item \textbf{enrollmentType} -- Tipo de trajeto (ONE\_WAY, RETURN, ROUND\_TRIP)
    \item \textbf{recordedPrice} -- Preço registrado no momento da inscrição (imutável)
    \item \textbf{userId} -- FK: Usuário
    \item \textbf{tripInstanceId} -- FK: Instância de Viagem
    \item \textbf{boardingStop} -- Parada de embarque
    \item \textbf{alightingStop} -- Parada de desembarque
    \item \textbf{organizationId} -- FK: Organização (multi-tenancy)
    \item \textbf{activeKey} -- Chave única para garantir no máximo uma inscrição ativa por (userId, tripInstanceId)
    \item \textbf{createdAt, updatedAt} -- Timestamps
\end{itemize}

\paragraph{\textbf{Payment (Pagamentos)}}\mbox{}\\
Registra transações de pagamento associadas a inscrições (relação 1:1 com Booking).
\begin{itemize}
    \item \textbf{id} -- Identificador único (UUID)
    \item \textbf{organizationId} -- FK: Organização
    \item \textbf{method} -- Método de pagamento (MONEY, PIX, CREDIT\_CARD, DEBIT\_CARD)
    \item \textbf{amount} -- Valor (Decimal 10.2)
    \item \textbf{status} -- Status (PENDING, COMPLETED, FAILED)
    \item \textbf{enrollmentId} -- FK: Inscrição (\textit{unique}, 1:1 com Booking)
    \item \textbf{createdAt, updatedAt} -- Timestamps
\end{itemize}

\paragraph{\textbf{TripSchedulingConfig (Configuração de Agendamento)}}\mbox{}\\
Configuração por organização (1:0..1) que controla os \textit{jobs} de geração recorrente e auto-cancelamento de viagens.
\begin{itemize}
    \item \textbf{id} -- Identificador único (UUID)
    \item \textbf{organizationId} -- FK: Organização (único)
    \item \textbf{daysAhead} -- Janela de geração antecipada de instâncias (1--90, padrão 14)
    \item \textbf{enabled} -- Se os \textit{jobs} estão ativos para a organização
    \item \textbf{createdAt, updatedAt} -- Timestamps
\end{itemize}

\paragraph{\textbf{AuditLog (Log de Auditoria)}}\mbox{}\\
Registra todas as ações críticas realizadas no sistema para conformidade e rastreabilidade.
\begin{itemize}
    \item \textbf{id} -- Identificador único (UUID)
    \item \textbf{organizationId} -- FK: Organização
    \item \textbf{userId} -- FK: Usuário que realizou a ação
    \item \textbf{action} -- Descrição da ação realizada
    \item \textbf{timestamp} -- Data e hora da ação
\end{itemize}

\subsection{Interfaces}
As interfaces de usuário (telas e protótipos da aplicação cliente) são apresentadas na Seção~\ref{sec:frontend}, dedicada ao \textit{frontend}.

% ------------------------------
\section{Implementação e Resultados Alcançados}

O desenvolvimento do projeto atingiu 100\% do escopo de módulos previsto (13 módulos completos), com implementação robusta de todos os componentes fundamentais para operação do sistema em produção. Todos os módulos seguem rigorosamente os padrões de Clean Architecture e são estruturados com casos de uso, value objects, repositórios e tratamento de erros específicos de domínio.

\subsection{Módulos Implementados (13 Módulos --- 100\% Completo)}

A API expõe 93 endpoints REST distribuídos pelos módulos de domínio, totalizando 93 casos de uso. A Tabela~\ref{tab:endpoints} consolida a contagem por módulo; nas subseções seguintes, cada módulo é descrito com seus casos de uso e endpoints representativos (a especificação completa das rotas está disponível via Swagger em \texttt{/api}).

\begin{table}[H]
\centering
\caption{Endpoints REST e casos de uso por módulo}
\label{tab:endpoints}
\begin{tabular}{lrr}
\toprule
\textbf{Módulo} & \textbf{Endpoints} & \textbf{Use Cases} \\
\midrule
Auth & 9 & 10 \\
User & 8 & 6 \\
Organization & 9 & 8 \\
Membership & 8 & 8 \\
Driver & 9 & 9 \\
Vehicle & 5 & 5 \\
Trip & 17 & 19 \\
Bookings & 10 & 10 \\
Plans & 6 & 6 \\
Subscriptions & 5 & 5 \\
Payment & 4 & 4 \\
Scheduling & 2 & 2 \\
Plan-Usage & 1 & 1 \\
\midrule
\textbf{Total} & \textbf{93} & \textbf{93} \\
\bottomrule
\end{tabular}
\end{table}

\subsubsection{Role Management}

\textbf{Descrição}: Sistema de papéis (roles) que define permissões de acesso (Status: Implementado).

\textbf{Funcionalidades}:
\begin{itemize}
\item Entidade Role com enum de dois papéis: ADMIN e DRIVER.
\item Repositório de Roles com persistência e recuperação.
\item Seed automático que popula roles no banco no primeiro boot.
\item Integração com Docker Compose para inicialização automática.
\end{itemize}

\textbf{Endpoints}: N/A (Roles são pré-definidas no sistema)

\subsubsection{Shared Module}

\textbf{Descrição}: Módulo centralizado que expõe componentes reutilizáveis em todos os contextos delimitados (Status: Implementado).

\textbf{Componentes Exportados}:
\begin{itemize}
\item PrismaModule: Conexão e injeção com o banco de dados.
\item JwtGuard: Guard global para autenticação JWT.
\item RolesGuard, TenantFilterGuard, DevGuard: Guards de autorização.
\item LoggingInterceptor: Interceptador de logging centrado.
\item AllExceptionsFilter: Filter global de tratamento de exceções.
\item Value Objects: Email, Telephone, Money (com validações).
\item UnitOfWork + DbContext: Transações atômicas via AsyncLocalStorage.
\item Domain Errors: Classes base para erros de negócio.
\end{itemize}

\subsubsection{User Module (CRUD + Autenticação)}

\textbf{Descrição}: Gerenciamento completo de usuários e autenticação via JWT (Status: Implementado).

\textbf{Use Cases Implementados}:
\begin{enumerate}
\item CreateUserUseCase: Cadastro com validação de email único e hash de senha Bcrypt.
\item FindAllUsersUseCase: Listagem paginada de usuários ativos.
\item FindUserByIdUseCase: Busca individual com validação 404.
\item UpdateUserUseCase: Atualização com re-validação de dados.
\item DeleteUserUseCase: Soft-delete alterando status para INACTIVE.
\item LoginUseCase: Autenticação com geração de JWT + refresh token.
\item RegisterUseCase: Auto-registro de novo usuário.
\item RefreshTokenUseCase: Renovação de token expirado.
\item RegisterOrganizationWithAdminUseCase: Fluxo unificado de onboarding — cria usuário, organização e membership ADMIN atomicamente via \texttt{UnitOfWork}, com rollback automático pelo Prisma em caso de falha.
\item SetupOrganizationForExistingUserUseCase: Cria organização para usuário já autenticado, gera membership ADMIN e re-emite JWT com o novo \texttt{organizationId} no payload.
\end{enumerate}

\textbf{Endpoints}: 9 rotas de autenticação --- \texttt{POST /auth/\{login, register, refresh, logout, register-organization, setup-organization, forgot-password, reset-password, verify-email\}} --- e o CRUD de usuários (\texttt{POST/GET/PUT/DELETE /users}, com \texttt{GET} paginado e \textit{soft-delete} no \texttt{DELETE}).

\textbf{Segurança}: Bcrypt (salt rounds: 10) para hash de senhas; JWT com secrets distintos para access (1h) e refresh tokens (7d). O payload do JWT é enriquecido no login com todos os dados necessários, eliminando consultas ao banco em cada requisição. O módulo implementa ainda \textbf{recuperação de senha} e \textbf{verificação de e-mail} por tokens de uso único armazenados apenas como hash SHA-256 (TTL de 1h e 24h), com resposta constante na solicitação (anti-enumeration) e revogação de \textit{refresh tokens} via \texttt{jti} no \textit{logout} e na redefinição de senha.

\subsubsection{Organization Module (CRUD)}

\textbf{Descrição}: Gerenciamento de organizações de transporte com multi-tenancy (Status: Implementado).

\textbf{Use Cases Implementados}:
\begin{enumerate}
\item CreateOrganizationUseCase: Criação com slug auto-gerado e validação CNPJ.
\item FindAllOrganizationsUseCase: Listagem paginada de todas (ativas e inativas).
\item FindAllActiveOrganizationsUseCase: Apenas organizações com status ACTIVE.
\item FindOrganizationByIdUseCase: Busca individual com erro 404 e \texttt{OrganizationForbiddenError} (HTTP 403).
\item UpdateOrganizationUseCase: Atualização com re-validação.
\item DisableOrganizationUseCase: Soft-delete com timestamp.
\item FindOrganizationByUserUseCase: Retorna orgs associadas ao \texttt{userId} do token JWT.
\end{enumerate}

\textbf{Endpoints}: 9 rotas. Exemplos: \texttt{GET /organizations/me} (orgs do usuário autenticado), \texttt{GET /organizations/active} (apenas ativas, paginado), \texttt{PUT/DELETE /organizations/:id} (atualização e \textit{soft-delete}). Há ainda endpoints públicos sem autenticação --- \texttt{GET /public/organizations} (lista paginada) e \texttt{GET /public/organizations/:slug} --- para páginas institucionais.

\textbf{Value Objects}: CNPJ (com dígitos verificadores), OrganizationName, Slug (URL-safe), Address.

\subsubsection{Membership Module}

\textbf{Descrição}: Associações entre usuários, roles e organizações. Base fundamental para RBAC (Status: Implementado).

\textbf{Estrutura}: Tabela de junção \textit{OrganizationMembership} com chave composta (userId, roleId, organizationId).

\textbf{Use Cases Implementados}:
\begin{enumerate}
\item CreateMembershipUseCase: O \texttt{organizationId} vem exclusivamente do JWT, eliminando o vetor de ataque cross-tenant. Valida pré-requisito de perfil Driver antes de criar membership com role DRIVER.
\item FindMembershipByCompositeKeyUseCase: Busca específica.
\item FindMembershipsByUserUseCase: Listagem paginada filtrada pela org do caller.
\item FindMembershipsByOrganizationUseCase: Listagem paginada por organização.
\item RemoveMembershipUseCase: Soft-delete via \texttt{removedAt} timestamp.
\item RestoreMembershipUseCase: Reversão de soft-delete.
\item FindRoleByUserIdAndOrganizationIdUseCase: Retorna o role do usuário em uma organização.
\end{enumerate}

\textbf{Endpoints}: 8 rotas operando sobre a chave composta. Exemplos: \texttt{POST /memberships} (org vem do JWT), \texttt{GET /memberships/me/role/:organizationId}, \texttt{DELETE} (\textit{soft-delete}) e \texttt{PATCH .../restore} sobre \texttt{/:userId/:roleId/:organizationId}.

\textbf{Erros de Domínio}: \texttt{MembershipAlreadyExistsError} (HTTP 409), \texttt{MembershipNotFoundError} (HTTP 404), \texttt{DriverNotFoundForMembershipError} (HTTP 400).

\subsubsection{Driver Module}

\textbf{Descrição}: Gerenciamento de perfis de motoristas com Value Objects para CNH. Motoristas são entidades globais — o vínculo com organizações ocorre exclusivamente via \texttt{OrganizationMembership}, permitindo que um usuário seja motorista em múltiplas organizações (Status: Implementado).

\textbf{Fluxo de Onboarding de Motorista}:
\begin{enumerate}
\item Usuário cria perfil via \texttt{POST /drivers} (self-service; \texttt{userId} do JWT).
\item Admin verifica identidade via \texttt{GET /drivers/lookup?email=\&cnh=}.
\item Admin vincula à organização via \texttt{POST /memberships} com \texttt{roleId=DRIVER}.
\end{enumerate}

\textbf{Use Cases Implementados (9 total)}:
\begin{enumerate}
\item CreateDriverUseCase: Criação self-service com check de duplicata (\texttt{DriverAlreadyExistsError} HTTP 409).
\item UpdateDriverUseCase: Atualização (fluxo admin) com verificação de ownership.
\item UpdateOwnDriverUseCase: Atualização parcial self-service (\texttt{cnhExpiresAt}/\texttt{cnhCategories}) via \texttt{PATCH /drivers/me}.
\item FindDriverByIdUseCase: Busca com proteção IDOR via \texttt{belongsToOrganization}.
\item FindDriverByUserIdUseCase: Busca por usuário.
\item FindDriverNameByIdUseCase: Retorna o nome do motorista (visível a passageiros).
\item FindAllDriversByOrganizationUseCase: Paginação com JOIN via membership.
\item RemoveDriverUseCase: Soft delete com validação de ownership.
\item LookupDriverUseCase: Verificação cruzada email + CNH para administrador.
\end{enumerate}

\textbf{Value Objects}: \textbf{Cnh} (9--12 caracteres alfanuméricos) e \textbf{CnhCategories} --- coleção de categorias (A--E) com deduplicação, ordenação e normalização \textit{case-insensitive}, permitindo que um motorista possua múltiplas categorias (ex.: A+B).

\textbf{Endpoints}: 9 rotas. Exemplos: \texttt{POST /drivers} (perfil \textit{self-service}), \texttt{GET /drivers/me}, \texttt{PATCH /drivers/me} (atualização parcial de \texttt{cnhExpiresAt}/\texttt{cnhCategories}), \texttt{GET /drivers/lookup?email=\&cnh=} (ADMIN), \texttt{GET /drivers/organization/:orgId} (paginado) e \texttt{GET /drivers/:id/name} (visível a qualquer usuário autenticado).

\subsubsection{Vehicle Module}

\textbf{Descrição}: Gerenciamento de frota de veículos com Value Object para placa e proteção IDOR (Status: Implementado).

\textbf{Use Cases Implementados (5 total)}:
\begin{enumerate}
\item CreateVehicleUseCase: Criação com validação de placa única (\texttt{PlateAlreadyInUseError} HTTP 409).
\item FindVehicleByIdUseCase: Busca com proteção IDOR via \texttt{organizationId} do JWT.
\item FindAllVehiclesByOrganizationUseCase: Listagem paginada por organização.
\item UpdateVehicleUseCase: Atualização com proteção IDOR e verificação de status ativo (\texttt{VehicleInactiveError} HTTP 410).
\item RemoveVehicleUseCase: Soft-delete com proteção IDOR.
\end{enumerate}

\textbf{Value Objects}: \textbf{Plate} — valida placa brasileira nos formatos \texttt{ABC1234} (padrão) e \texttt{ABC1D23} (Mercosul).

\textbf{Endpoints}: 5 rotas. Exemplos: \texttt{POST/GET /vehicles/organization/:organizationId} (registro e listagem paginada) e \texttt{GET/PUT/DELETE /vehicles/:id} (consulta, atualização e \textit{soft-delete}, todas com proteção IDOR).

\subsubsection{Trip Module}

\textbf{Descrição}: Módulo de viagens com separação entre templates (configuração recorrente) e instâncias (execuções reais), máquina de estados e atribuição de motorista/veículo (Status: Implementado).

\textbf{TripTemplate --- Use Cases (5 total)}:
\begin{enumerate}
\item CreateTripTemplateUseCase: Cria template com configuração de preços por tipo de trajeto.
\item FindTripTemplateByIdUseCase: Busca com proteção IDOR.
\item FindAllTripTemplatesByOrganizationUseCase: Listagem paginada por organização.
\item UpdateTripTemplateUseCase: Atualização com validações de acesso e status.
\item DeactivateTripTemplateUseCase: Soft-deactivation via \texttt{status = INACTIVE}.
\end{enumerate}

\textbf{TripInstance --- Use Cases (7 total)}:
\begin{enumerate}
\item CreateTripInstanceUseCase: Criação a partir de template com propagação de campos.
\item FindTripInstanceByIdUseCase: Busca individual com proteção IDOR.
\item FindAllTripInstancesByOrganizationUseCase: Listagem paginada por organização.
\item FindTripInstancesByTemplateUseCase: Listagem paginada por template.
\item TransitionTripInstanceStatusUseCase: Máquina de estados da viagem.
\item AssignDriverToTripInstanceUseCase: Atribuição/desatribuição de motorista com validação de existência.
\item AssignVehicleToTripInstanceUseCase: Atribuição/desatribuição de veículo com validação de existência.
\end{enumerate}

\textbf{Endpoints}: 17 rotas. \textit{TripTemplate}: CRUD em \texttt{/trip-templates} mais \texttt{POST /trip-templates/:id/generate-instances} (geração manual). \textit{TripInstance}: \texttt{POST /trip-instances}, listagens por org/template, \texttt{GET /trip-instances/driver/me} (viagens do motorista), \texttt{PATCH /trip-instances/:id/status} (máquina de estados) e \texttt{:id/driver}, \texttt{:id/vehicle} (atribuição). Há ainda rotas públicas \texttt{GET /public/trip-instances} e \texttt{.../org/:slug}.

\textbf{Máquina de Estados}:

\begin{table}[H]
\centering
\caption{Transições de status de TripInstance}
\label{tab:trip_status}
\begin{tabular}{lccccc}
\toprule
\textbf{De $\backslash$ Para} & \textbf{SCHEDULED} & \textbf{CONFIRMED} & \textbf{IN\_PROGRESS} & \textbf{FINISHED} & \textbf{CANCELED} \\
\midrule
DRAFT & \checkmark* & -- & -- & -- & \checkmark \\
SCHEDULED & -- & \checkmark* & -- & -- & \checkmark \\
CONFIRMED & \checkmark & -- & \checkmark* & -- & \checkmark \\
IN\_PROGRESS & -- & -- & -- & \checkmark & \checkmark \\
FINISHED & -- & -- & -- & -- & -- \\
CANCELED & -- & -- & -- & -- & -- \\
\bottomrule
\end{tabular}
\end{table}
\footnotesize{* Exige driver e vehicle atribuídos. Transições inválidas lançam \texttt{InvalidTripStatusTransitionError} (HTTP 422).}

\subsubsection{Bookings Module}

\textbf{Descrição}: Gerenciamento completo de inscrições de passageiros em viagens, com controle de capacidade, preço registrado server-side e proteção contra race conditions (Status: Implementado).

\textbf{Use Cases Implementados (9 total)}:
\begin{enumerate}
\item CreateBookingUseCase: Inscreve passageiro com validação de status da viagem, controle de capacidade e preço calculado server-side a partir do template. Todo o fluxo é executado atomicamente via \texttt{UnitOfWork} com isolamento \texttt{Serializable}, prevenindo race conditions clássicas de sistemas de reservas.
\item CancelBookingUseCase: Cancela inscrição com bloqueio para viagens \texttt{IN\_PROGRESS}/\texttt{FINISHED} e prazo de 30 minutos antes da partida.
\item ConfirmPresenceUseCase: Confirma presença do passageiro (apenas org members; owner bloqueado).
\item FindBookingByIdUseCase: Busca com isolamento de tenant.
\item FindBookingsByOrganizationUseCase: Listagem paginada por organização.
\item FindBookingsByTripInstanceUseCase: Lista passageiros de uma viagem (apenas org members).
\item FindBookingsByUserUseCase: Bookings do próprio usuário com filtro de status opcional.
\item FindBookingDetailsUseCase: Detalhe enriquecido com dados da \texttt{TripInstance} (horário, capacidade, vagas disponíveis).
\item GetBookingAvailabilityUseCase: Verifica disponibilidade de vagas antes da inscrição.
\end{enumerate}

\textbf{Endpoints}: 10 rotas. Exemplos: \texttt{POST /bookings} (preço \textit{server-side}), \texttt{GET /bookings/availability/:tripInstanceId} (vagas), \texttt{GET /bookings/:id/details} (detalhe enriquecido), \texttt{PATCH /bookings/:id/cancel} e \texttt{.../confirm-presence} (confirmação restrita a \textit{org members}).

\textbf{Erros de Domínio}: \texttt{BookingNotFoundError} (404), \texttt{BookingAlreadyExistsError} (409), \texttt{TripInstanceFullError} (409), \texttt{TripInstanceNotBookableError} (400), \texttt{BookingCancellationNotAllowedError} (400), \texttt{BookingCancellationDeadlineError} (400), \texttt{BookingAlreadyInactiveError} (400), \texttt{BookingAccessForbiddenError} (403), entre outros.

\subsubsection{Plans Module}

\textbf{Descrição}: Gerenciamento dos planos de assinatura disponíveis na plataforma. Escritas restritas a desenvolvedores via \texttt{DevGuard} (Status: Implementado).

\textbf{Use Cases (5 total)}: \texttt{CreatePlanUseCase}, \texttt{UpdatePlanUseCase}, \texttt{DeactivatePlanUseCase}, \texttt{FindPlanByIdUseCase}, \texttt{FindAllPlansUseCase}.

\textbf{Endpoints}: 6 rotas. Escritas restritas a \texttt{@Dev()} (\texttt{POST /plans}, \texttt{PATCH /plans/:id}, \texttt{.../deactivate}); leitura autenticada (\texttt{GET /plans}, \texttt{/plans/:id}) e endpoint público \texttt{GET /public/plans} para a tela de \textit{onboarding}.

\subsubsection{Subscriptions Module}

\textbf{Descrição}: Gerenciamento de assinaturas das organizações nos planos. Implementa expiração lazy (sem cron job) e controle de limites de plano via \texttt{PlanLimitService} (Status: Implementado).

\textbf{Expiração Lazy}: O helper \texttt{resolveActiveSubscription(orgId, repo)} centraliza a busca + verificação de expiração. Se a assinatura estiver vencida, é transicionada para \texttt{PAST\_DUE} on-demand no momento da leitura.

\textbf{PlanLimitService}: Centraliza verificação de limites de plano (veículos, motoristas, viagens/mês). Injetável em qualquer use case que precise bloquear criação de recursos além da cota.

\textbf{Use Cases (4 total)}:
\begin{enumerate}
\item SubscribeToPlanUseCase: Valida plano ativo e ausência de assinatura ACTIVE. Calcula \texttt{expiresAt = startDate + plan.durationDays}.
\item CancelSubscriptionUseCase: Transita para CANCELED.
\item FindActiveSubscriptionUseCase: Retorna assinatura ativa com lazy expiry.
\item FindSubscriptionsByOrganizationUseCase: Histórico paginado.
\end{enumerate}

\textbf{Endpoints}: 5 rotas sob \texttt{/organizations/:orgId/subscriptions}. Exemplos: \texttt{POST} (assinar plano), \texttt{PATCH /:id/cancel}, \texttt{GET /active} (assinatura vigente, com \textit{lazy expiry}) e \texttt{GET} (histórico paginado).

\subsubsection{Payment Module}

\textbf{Descrição}: Registro e simulação de pagamentos associados a inscrições. A criação de pagamentos é responsabilidade do \texttt{CreateBookingUseCase} — cada inscrição gera automaticamente um registro de pagamento com status \texttt{PENDING}. Os endpoints de simulação permitem demonstrar o fluxo completo sem gateway externo (Status: Implementado). Além do ADMIN, o motorista pode confirmar/falhar pagamentos de viagens atribuídas a ele (\texttt{PaymentNotAssignedToDriverError}, HTTP 403). A resposta inclui \texttt{tripInstanceId} e \texttt{tripDepartureTime} (\textit{snapshots} de leitura), permitindo ao cliente agregar receita por data da viagem.

\textbf{Use Cases (4 total)}: \texttt{FindPaymentByIdUseCase}, \texttt{FindPaymentsByOrganizationUseCase}, \texttt{ConfirmPaymentUseCase} (\texttt{PENDING → COMPLETED}), \texttt{FailPaymentUseCase} (\texttt{PENDING → FAILED}).

\textbf{Endpoints}: 4 rotas sob \texttt{/organizations/:orgId/payments}. \texttt{GET /:id} e \texttt{GET} (listagem paginada); \texttt{PATCH /:id/confirm} (\texttt{PENDING → COMPLETED}) e \texttt{.../fail} (\texttt{PENDING → FAILED}), simulando o ciclo sem gateway externo.

\subsubsection{Scheduling Module}

\textbf{Descrição}: Automação da geração e do encerramento de viagens recorrentes via \textit{jobs} agendados (\texttt{@nestjs/schedule}), parametrizados por organização através de \texttt{TripSchedulingConfig} (Status: Implementado).

\textbf{Jobs (\textit{cron})}:
\begin{enumerate}
\item \textbf{Geração recorrente} (\texttt{0 2 * * *} UTC): para cada organização ativa, materializa uma janela rolante de \texttt{TripInstance} a partir dos \texttt{TripTemplate} recorrentes, respeitando \texttt{daysAhead} e a cota de viagens do plano. Idempotente por (template, dia).
\item \textbf{Auto-cancelamento} (\texttt{*/15 * * * *} UTC): transiciona para \texttt{CANCELED} as instâncias cujo prazo (\texttt{autoCancelAt}) venceu sem atingir a receita mínima, propagando a falha aos pagamentos \texttt{PENDING} associados.
\end{enumerate}

\textbf{Use Cases (3 total)}: \texttt{FindTripSchedulingConfigUseCase}, \texttt{UpdateTripSchedulingConfigUseCase} e a geração manual \texttt{GenerateTripInstancesForTemplateUseCase} (mesmo núcleo do \textit{cron}, exposta via \texttt{POST /trip-templates/:id/generate-instances}).

\textbf{Endpoints}: 2 rotas sob \texttt{/organizations/:organizationId/scheduling-config} (\texttt{GET} e \texttt{PATCH}, ADMIN). A cadência dos \textit{jobs} é fixa; a configuração por organização controla apenas a janela (\texttt{daysAhead}) e a ativação (\texttt{enabled}). O \texttt{ScheduleModule} é registrado condicionalmente (\texttt{DISABLE\_CRON}).

\subsubsection{Plan-Usage Module}

\textbf{Descrição}: Exposição consolidada do consumo de recursos da organização frente aos limites do plano vigente, para alimentar painéis administrativos (Status: Implementado).

\textbf{Use Case (1 total)}: \texttt{FindPlanUsageUseCase} --- agrega contagem de veículos, motoristas e viagens (estas contadas dentro do \textbf{período de assinatura} corrente, \texttt{[startDate, expiresAt)}, e não por mês-calendário) e confronta com \texttt{maxVehicles}, \texttt{maxDrivers} e \texttt{maxMonthlyTrips}.

\textbf{Endpoints}: 1 rota --- \texttt{GET /organizations/:id/plan-usage}.

\subsection{Infraestrutura Transacional --- UnitOfWork + DbContext}

Para garantir atomicidade em operações que envolvem múltiplas escritas (ex: criar \texttt{Booking} + \texttt{Payment}), foi implementada uma infraestrutura de transações baseada em \texttt{AsyncLocalStorage}:

\begin{lstlisting}[caption={Fluxo de transação via UnitOfWork}, label={cod:uow}]
UseCase
  -> unitOfWork.execute(async () => {
       await bookingRepository.save(booking)
       await paymentRepository.save(payment)
     })
         |
PrismaUnitOfWork
  -> prisma.$transaction((tx) =>
       dbContext.run(tx, fn))   # coloca tx no AsyncLocalStorage
           |
DbContext (AsyncLocalStorage)
  -> client = tx durante a execucao
           |
Repository -> usa dbContext.client automaticamente
\end{lstlisting}

Fora de transação, \texttt{dbContext.client} aponta para o \texttt{PrismaService} (cliente raiz). Dentro da transação, aponta para o \texttt{tx}. A camada de aplicação não conhece Prisma e não precisa alterar assinaturas de repositórios.

\subsection{Sistema de Autenticação e Autorização}

\subsubsection{Fluxo de Autenticação JWT}

\begin{enumerate}
\item \textbf{Registro}: Password hasheado com Bcrypt (10 salt rounds) antes de persistência.
\item \textbf{Login}: Sistema valida credenciais e gera dois tokens:
    \begin{itemize}
    \item \textbf{Access Token}: JWT com expiração curta (1 hora).
    \item \textbf{Refresh Token}: JWT com expiração longa (7 dias).
    \end{itemize}
\item \textbf{Refresh de Token}: Cliente envia refresh token; sistema emite novo par.
\item \textbf{Proteção de Rotas}: \texttt{JwtAuthGuard} valida access token. Token inválido ou expirado resulta em 401.
\end{enumerate}

\subsubsection{RBAC (Role-Based Access Control)}

\begin{itemize}
\item \textbf{ADMIN}: Acesso total. Pode gerenciar organizações, usuários, motoristas e configurações.
\item \textbf{DRIVER}: Acesso limitado. Pode visualizar viagens e confirmar ações relacionadas.
\item \textbf{USER} (Implícito): Passageiros. Podem se inscrever em viagens e visualizar suas inscrições.
\end{itemize}

Memberships permitem múltiplas roles por usuário em diferentes organizações (ex: motorista em uma org, admin em outra).

\subsection{Multi-tenancy}

O isolamento lógico de dados entre organizações é implementado através de:
\begin{enumerate}
\item \textbf{Campo organizationId}: Presente em todas as entidades que pertencem a uma organização.
\item \textbf{Filtragem em Queries}: Todas as buscas filtram por \texttt{organizationId} da organização do usuário autenticado.
\item \textbf{TenantFilterGuard}: Valida que o parâmetro \texttt{:organizationId} na rota corresponde ao \texttt{organizationId} do JWT.
\item \textbf{Proteção IDOR (OWASP A01)}: Endpoints com parâmetro \texttt{/:id} validam ownership via \texttt{organizationId} do JWT em Driver, Vehicle e Trip.
\end{enumerate}

\subsection{Tratamento de Erros Estruturado}

O sistema implementa hierarquia de exceções mapeadas pelo \texttt{AllExceptionsFilter} global:

\begin{table}[H]
\centering
\caption{Mapeamento de erros de domínio para HTTP}
\label{tab:error_mapping}
\begin{tabular}{ll}
\toprule
\textbf{Sufixo do código de erro} & \textbf{HTTP} \\
\midrule
\texttt{\_NOT\_FOUND} & 404 Not Found \\
\texttt{\_ALREADY\_EXISTS} & 409 Conflict \\
\texttt{INVALID\_} / \texttt{\_BAD\_REQUEST} & 400 Bad Request \\
\texttt{\_FORBIDDEN} & 403 Forbidden \\
\texttt{\_UNAUTHORIZED} & 401 Unauthorized \\
\texttt{\_GONE} & 410 Gone \\
\bottomrule
\end{tabular}
\end{table}

\subsection{Banco de Dados: Migrações e Seed}

Prisma gerencia histórico completo de mudanças no esquema com migrations versionadas. O script \texttt{prisma/seed.ts} popula automaticamente roles padrão (ADMIN, DRIVER), planos (FREE, BASIC, PRO, PREMIUM) e um usuário admin de teste.

\subsection{Cobertura de Testes}

\textbf{Status Atual}: 458 testes em 57 suites, cobrindo os use cases críticos de todos os módulos, guards de autorização e fluxos de negócio complexos.

\textbf{Padrão de Testes Adotado}:
\begin{itemize}
\item \textbf{AAA (Arrange-Act-Assert)}: Estrutura consistente em todas as suites.
\item \textbf{makeMocks()}: Cria todos os mocks via \texttt{jest.fn()} com tipagem \texttt{jest.Mocked<T>}.
\item \textbf{setupHappyPath()}: Configura mocks para o cenário de sucesso padrão.
\item \textbf{Injeção manual}: Dependências injetadas diretamente, sem mocks de framework NestJS.
\item \textbf{Factories}: Funções \texttt{make*()} por módulo com suporte a overrides parciais.
\end{itemize}

\textbf{Distribuição das suites por módulo (57 suites, 458 testes)}: a Tabela~\ref{tab:testes} resume a cobertura agregada por módulo; o detalhamento por use case está versionado no repositório (\texttt{test/modules/}).

\begin{table}[H]
\centering
\caption{Cobertura de testes unitários agregada por módulo}
\label{tab:testes}
\begin{tabular}{lrr}
\toprule
\textbf{Módulo} & \textbf{Suites} & \textbf{Testes} \\
\midrule
Trip & 18 & 183 \\
Bookings & 9 & 88 \\
Auth & 7 & 36 \\
Driver & 4 & 33 \\
Shared (Guards) & 3 & 26 \\
Payment & 3 & 20 \\
Subscriptions & 3 & 18 \\
Scheduling & 2 & 15 \\
Membership & 1 & 10 \\
Organization & 2 & 10 \\
Vehicle & 2 & 7 \\
User & 2 & 6 \\
Plans & 1 & 6 \\
\midrule
\textbf{Total} & \textbf{57} & \textbf{458} \\
\bottomrule
\end{tabular}
\end{table}

Os testes apoiam-se em \textit{factories} por módulo (funções \texttt{make*()} com suporte a \textit{overrides} parciais) que constroem entidades de domínio e DTOs de entrada --- por exemplo \texttt{makeUser}, \texttt{makeOrganization}, \texttt{makeDriver}, \texttt{makeTripTemplate}, \texttt{makeTripInstance}, \texttt{makeBooking}, \texttt{makeSubscription} ---, eliminando duplicação de \textit{setup} entre as suites.

% =========================================================
% SEÇÃO FRONTEND --- PLACEHOLDER ESTRUTURADO
% TODO: preencher com o conteúdo da aplicação cliente.
% Esta seção foi pré-estruturada para a integração da documentação
% do frontend ao TCC. Remover este comentário após o preenchimento.
% =========================================================
\section{Aplicação Cliente (Frontend)}
\label{sec:frontend}

% TODO: parágrafo introdutório — propósito da aplicação cliente, público-alvo
% (estudantes/passageiros, administradores, motoristas) e como ela consome a Movy API.

\subsection{Tecnologias do Frontend}

% TODO: descrever a stack do frontend (framework/biblioteca, linguagem, gerenciamento
% de estado, estilização, build, etc.) e justificar as escolhas, no mesmo formato da
% Seção "Tecnologias Implementadas" do backend.
% Sugestão: replicar a tabela abaixo preenchendo as linhas.
%
% \begin{table}[H]
% \centering
% \caption{Tecnologias implementadas no frontend}
% \label{tab:tecnologias_frontend}
% \begin{tabular}{lll}
% \toprule
% \textbf{Tecnologia} & \textbf{Versão} & \textbf{Descrição} \\
% \midrule
% ... & ... & ... \\
% \bottomrule
% \end{tabular}
% \end{table}

\textit{[Conteúdo a ser inserido: stack tecnológica do frontend e justificativas.]}

\subsection{Arquitetura e Organização}

% TODO: organização de pastas/camadas, padrões adotados (ex.: componentização,
% camadas de serviços para chamadas HTTP, gerência de autenticação/refresh token).

\textit{[Conteúdo a ser inserido.]}

\subsection{Telas e Funcionalidades}

% TODO: descrever as principais telas e inserir prints/protótipos
% (login, onboarding de organização, listagem de viagens, inscrição, painel admin, etc.).
% Usar \begin{figure}[H] ... \includegraphics ... \end{figure} para cada captura.

\textit{[Conteúdo a ser inserido: capturas de tela e descrição dos fluxos.]}

\subsection{Integração com a API}

% TODO: descrever como o frontend consome a Movy API — autenticação JWT (access/refresh),
% tratamento de erros padronizados (mapeamento de códigos HTTP), consumo de endpoints
% públicos (/public/organizations, /public/plans, /public/trip-instances) e fluxo
% multi-tenant (seleção/contexto de organização).

\textit{[Conteúdo a ser inserido.]}

% ------------------------------
\section{Considerações Finais}

\subsection{Síntese dos Resultados}

As metas estabelecidas para o projeto foram alcançadas com sucesso, atingindo 100\% do escopo de módulos planejado (13 módulos). O sistema demonstra uma arquitetura sólida baseada em Clean Architecture e Domain-Driven Design, funcional e pronta para uso em produção.

\subsubsection{Sucessos Implementados}

\begin{enumerate}
\item \textbf{Arquitetura Modular}: Cada módulo segue padrão consistente (Application → Domain → Infrastructure → Presentation), facilitando manutenção e testes.

\item \textbf{Multi-tenancy}: Implementado com sucesso em nível de banco de dados (\texttt{organizationId} presente em todas as entidades), permitindo que múltiplas organizações usem a plataforma com isolamento lógico de dados.

\item \textbf{Autenticação Robusta}: Sistema JWT com access tokens de curta duração e refresh tokens, implementado com Bcrypt (10 salt rounds). JWT Strategy otimizada --- sem consulta ao banco por request.

\item \textbf{RBAC Completo}: Guards (RolesGuard, TenantFilterGuard, DevGuard) e decorators (@Roles, @Dev) implementados, com suporte a múltiplas roles por usuário em diferentes organizações.

\item \textbf{Value Objects}: Encapsulamento de validações críticas (CNPJ, Email, Slug, CNH, Placa brasileira/Mercosul, Money) em classes de domínio imutáveis.

\item \textbf{Soft-delete}: Implementado em todas as entidades críticas, preservando histórico de dados e integridade referencial.

\item \textbf{Segurança IDOR (OWASP A01) e Rate Limiting}: Endpoints com parâmetro \texttt{/:id} validam ownership via \texttt{organizationId} do JWT (Driver, Vehicle, Trip); proteção global de 60 req/min via \texttt{@nestjs/throttler}.

\item \textbf{Módulo de Viagens e Automação}: TripTemplate + TripInstance com máquina de estados completa, atribuição de motorista/veículo com validação de FK, e \textit{jobs} de geração recorrente e auto-cancelamento por organização.

\item \textbf{Módulo de Inscrições}: Bookings com controle de capacidade, preço calculado server-side, confirmação de presença e proteção contra race conditions via \texttt{UnitOfWork} com isolamento \texttt{Serializable}.

\item \textbf{Planos e Assinaturas}: Modelo SaaS com tiers de plano, limites de recursos (\texttt{PlanLimitService}), expiração lazy e histórico de assinaturas.

\item \textbf{Pagamentos}: Registro automático de pagamentos no momento da inscrição com simulação de ciclo de vida (PENDING → COMPLETED/FAILED).

\item \textbf{Transações Atômicas}: UnitOfWork + DbContext via AsyncLocalStorage garantem atomicidade em operações multi-repositório sem expor Prisma à camada de aplicação.

\item \textbf{Testes Unitários}: 458 testes em 57 suites com padrão AAA, factories por módulo e injeção manual de dependências.

\item \textbf{CI/CD}: GitHub Actions configurado para build, lint e testes automatizados a cada push.

\item \textbf{Tratamento de Erros Centralizado}: AllExceptionsFilter global com mapeamento de erros de domínio por padrão de código; \texttt{docker-compose} pronto com auto-seed e migrations automáticas.
\end{enumerate}

\subsubsection{Métricas do Projeto}

\begin{table}[H]
\centering
\caption{Métricas de Progresso da Movy API}
\label{tab:metricas_projeto}
\begin{tabular}{lrr}
\toprule
\textbf{Métrica} & \textbf{Total} & \textbf{Completo} \\
\midrule
Módulos & 13 & 13 (100\%) \\
Use Cases & 93 & 93 (100\%) \\
Endpoints REST & 93 & 93 (100\%) \\
Entidades de Domínio & 14 & 14 (100\%) \\
Migrações Executadas & 28 & 28 (100\%) \\
Cobertura de Testes (Use Cases críticos) & 57 suites & 458 testes \\
CI/CD (GitHub Actions) & Configurado & \checkmark \\
Rate Limiting & Configurado & 60 req/min \\
\bottomrule
\end{tabular}
\end{table}

\subsection{Limitações Identificadas}

\begin{enumerate}
\item \textbf{Integração com Gateway de Pagamento}: O módulo de pagamentos opera com simulação de ciclo de vida (PENDING → COMPLETED/FAILED). A integração com um gateway real (ex: Stripe) requer configuração externa e é a próxima evolução natural do módulo.

\item \textbf{Notificações Pendentes}: Sistema de notificações por email/SMS ainda não foi implementado, afetando requisitos RF17 e RF18.

\item \textbf{Cobertura de Testes Parcial}: Embora 458 testes cubram os use cases críticos e os guards principais, alguns use cases secundários ainda não possuem specs dedicadas. A meta de 80\% de cobertura completa está em andamento.

\item \textbf{Caching}: Não há implementação de caching (Redis) para otimização de consultas frequentes.

\item \textbf{Relatórios}: Funcionalidades de geração de relatórios e dashboards (RF15, RF16) ainda não foram desenvolvidas.
\end{enumerate}

\subsection{Lições Aprendidas}

\begin{enumerate}
\item \textbf{Clean Architecture é Essencial}: A separação clara entre camadas facilitou desenvolvimento iterativo e testes de componentes isolados.

\item \textbf{DDD Agregou Valor}: Value Objects bem definidos preveniram bugs de validação e documentam regras de negócio naturalmente.

\item \textbf{Multi-tenancy Desde o Início}: Implementar multi-tenancy na arquitetura inicial foi crítico; retroajustar seria custoso.

\item \textbf{Segurança IDOR Requer Vigilância Ativa}: Endpoints com parâmetro \texttt{/:id} sem \texttt{:organizationId} na rota exigem validação explícita de ownership via \texttt{organizationId} do JWT ou método \texttt{belongsToOrganization()} no repositório.

\item \textbf{Desacoplamento de Módulos Previne Ciclos}: O padrão Orchestrator (Auth → coordena User, Organization, Membership) eliminou dependências circulares e violações do SRP.

\item \textbf{Testes AAA com Injeção Manual são Superiores}: Dependências injetadas manualmente (sem framework) tornam os testes mais rápidos e isolados, e forçam o código a ser testável por design.

\item \textbf{UnitOfWork Abstrai Transações Corretamente}: Usar AsyncLocalStorage para propagar o cliente transacional evita expor Prisma à camada de aplicação sem alterar as assinaturas dos repositórios.
\end{enumerate}

\subsection{Recomendações para Continuação}

\subsubsection{Ações Imediatas Recomendadas}

\begin{enumerate}
\item \textbf{Integração de Pagamentos}: Configurar gateway real (Stripe) para substituir a simulação atual, conectando-o ao fluxo de inscrições.

\item \textbf{Expandir Cobertura de Testes}: Adicionar specs para use cases secundários de Driver, Vehicle e demais módulos. Meta: 80\% de cobertura.

\item \textbf{Notificações}: Implementar sistema de notificações via email/SMS (RF17, RF18) utilizando Bull/BullMQ para processamento assíncrono.

\item \textbf{Performance Profiling}: Usar ferramentas como Node.js Clinic.js para identificar gargalos antes de escalar.

\item \textbf{Backup Strategy}: Configurar backups automáticos do banco PostgreSQL em desenvolvimento e produção.
\end{enumerate}

\subsubsection{Tecnologias Futuras a Considerar}

\begin{itemize}
\item \textbf{Redis}: Para caching de dados frequentes e gerenciamento de sessões.
\item \textbf{Bull/BullMQ}: Para fila de jobs assíncronos (notificações, relatórios).
\item \textbf{Elasticsearch}: Para busca full-text em viagens e históricos.
\item \textbf{Sentry/Rollbar}: Para monitoramento de erros em produção.
\item \textbf{DataDog/New Relic}: Para APM (Application Performance Monitoring).
\end{itemize}

\subsection{Conclusão}

A Movy API demonstra viabilidade técnica e arquitetural para funcionar como uma plataforma SaaS de gerenciamento de transporte universitário. Os princípios de Clean Architecture e DDD forneceram fundação sólida que antecipa crescimento futuro sem sacrificar manutenibilidade.

Com 100\% dos módulos implementados — autenticação, usuários, organizações, roles, memberships, motoristas, veículos, viagens, inscrições de passageiros, planos, assinaturas e pagamentos — e um sistema robusto de RBAC, multi-tenancy, proteção IDOR e 458 testes unitários operacionais, o projeto entrega uma plataforma funcional ponta a ponta para o gerenciamento do transporte universitário intermunicipal.

O código segue padrões estabelecidos de engenharia de software, é facilmente integrável com aplicações front-end (web/mobile) via API REST e as decisões técnicas (NestJS, Prisma, PostgreSQL, Docker) garantem escalabilidade horizontal e suportam modelos de deployment modernos.

% ------------------------------
\newpage
\renewcommand{\refname}{}
\section*{Referências}
\bibliographystyle{abntex2-num}
\bibliography{referencias}

% ------------------------------
\newpage
\section*{Anexos}

\subsection*{Declaração de Entrega do Código-Fonte}
Declaro que o código-fonte foi entregue à coordenação do curso xxxxx

Assinatura: \rule{10cm}{0.4pt}

Data: \rule{1cm}{0.4pt}/\rule{1cm}{0.4pt}/\rule{2cm}{0.4pt}

\subsection*{Declaração de Submissão ao NIT}
Declaro que o software foi submetido ao NIT – IFPI xxxxx

Assinatura: \rule{10cm}{0.4pt}

Data: \rule{1cm}{0.4pt}/\rule{1cm}{0.4pt}/\rule{2cm}{0.4pt}

\end{document}
